\chapter{Maxwell Relations}

\section*{Introduction}

The Maxwell relations are a set of four equations in classical thermodynamics that connect partial derivatives of thermodynamic state functions. Derived by James Clerk Maxwell in his 1865 paper ``A Dynamical Theory of the Electromagnetic Field,'' they arise from the mathematical property that exact differentials of state functions are independent of the path of integration---equivalently, that mixed second partial derivatives commute. These relations allow thermodynamic quantities that are difficult to measure directly (such as the entropy change with volume) to be expressed in terms of quantities that are easily measured in the laboratory (such as the pressure change with temperature).

Thermodynamic potentials $U$ (internal energy), $H$ (enthalpy), $F$ (Helmholtz free energy), and $G$ (Gibbs free energy) are state functions of a system. Each has a natural set of independent variables, and the exactness of their differentials yields a Maxwell relation for each. For example, $(\partial T/\partial V)_{S} = -(\partial P/\partial S)_{V}$ tells us that the rate of temperature change with volume at constant entropy equals the negative of the rate of pressure change with entropy at constant volume. While the right-hand side involves entropy (hard to control directly), the left-hand side involves pressure and temperature (easy to measure), making these relations powerful tools for converting theoretical expressions into experimentally accessible forms.


The four Maxwell relations, one for each thermodynamic potential, are:

\begin{align}
\left(\frac{\partial T}{\partial V}\right)_{S} &= -\left(\frac{\partial P}{\partial S}\right)_{V}
\label{eq:maxwell1} \\
\left(\frac{\partial T}{\partial P}\right)_{S} &= \left(\frac{\partial V}{\partial S}\right)_{P}
\label{eq:maxwell2} \\
\left(\frac{\partial S}{\partial V}\right)_{T} &= \left(\frac{\partial P}{\partial T}\right)_{V}
\label{eq:maxwell3} \\
\left(\frac{\partial S}{\partial P}\right)_{T} &= -\left(\frac{\partial V}{\partial T}\right)_{P}
\label{eq:maxwell4}
\end{align}

\section*{Fundamental Equations Used}

The derivation starts from the exact differential of internal energy, $dU = T\,dS - P\,dV$, from the first law of thermodynamics.

\section*{Assumptions}

\textbf{Assumptions:} [To be filled.]

\section*{Mathematical Derivation}

\textbf{Step 1: Start from the exact differential of the internal energy.}

The first law of thermodynamics states that the internal energy $U$ is a state function with natural variables $S$ (entropy) and $V$ (volume). Its exact differential is:
\begin{equation}
dU = T\,dS - P\,dV.
\end{equation}
Since $U = U(S,V)$ is a smooth function of two variables, the mixed second partial derivatives must be equal (Schwarz's theorem):
\begin{equation}
\frac{\partial^{2}U}{\partial S\,\partial V} = \frac{\partial^{2}U}{\partial V\,\partial S}.
\end{equation}

\textbf{Step 2: Identify the first Maxwell relation.}

From $dU = T\,dS - P\,dV$, we read off:
\begin{equation}
\left(\frac{\partial U}{\partial S}\right)_{V} = T, \qquad \left(\frac{\partial U}{\partial V}\right)_{S} = -P.
\end{equation}
Differentiating the first with respect to $V$ and the second with respect to $S$:
\begin{equation}
\left(\frac{\partial T}{\partial V}\right)_{S} = -\left(\frac{\partial P}{\partial S}\right)_{V}.
\end{equation}

\textbf{Step 3: Derive the relation from enthalpy.}

The enthalpy is $H = U + PV$, a state function of $S$ and $P$:
\begin{equation}
dH = T\,dS + V\,dP.
\end{equation}
The same Schwarz condition gives:
\begin{equation}
\left(\frac{\partial T}{\partial P}\right)_{S} = \left(\frac{\partial V}{\partial S}\right)_{P}.
\end{equation}

\textbf{Step 4: Derive the relation from Helmholtz free energy.}

The Helmholtz free energy is $F = U - TS$, a state function of $T$ and $V$:
\begin{equation}
dF = -S\,dT - P\,dV.
\end{equation}
The exactness condition yields:
\begin{equation}
\left(\frac{\partial S}{\partial V}\right)_{T} = \left(\frac{\partial P}{\partial T}\right)_{V}.
\end{equation}

\textbf{Step 5: Derive the relation from Gibbs free energy.}

The Gibbs free energy is $G = H - TS$, a state function of $T$ and $P$:
\begin{equation}
dG = -S\,dT + V\,dP.
\end{equation}
The Schwarz condition gives:
\begin{equation}
\left(\frac{\partial S}{\partial P}\right)_{T} = -\left(\frac{\partial V}{\partial T}\right)_{P}.
\end{equation}

\section*{Characteristics of the Equation}

\begin{itemize}
  \item \textbf{Universality:} The Maxwell relations hold for any substance, ideal or real, as long as the thermodynamic potential is well-defined.
  \item \textbf{Experimental utility:} They convert ``hard'' derivatives (involving $S$) into ``easy'' ones (involving $P$, $V$, $T$), which are directly measurable.
  \item \textbf{Connection to exactness:} They are equivalent to the integrability condition $d^{2}F = 0$ for any exact differential $dF$.
  \item \textbf{Symmetry:} The four relations form a symmetric set connected to the four thermodynamic potentials through Legendre transforms.
\end{itemize}
\section*{Solution Methods}

\begin{enumerate}
  \item \textbf{From exact differentials:} Start with, e.g., $dU = T\,dS - P\,dV$. The mixed partial derivative condition $\partial^{2}U/\partial S\,\partial V = \partial^{2}U/\partial V\,\partial S$ yields relation~\eqref{eq:maxwell1}.
  \item \textbf{From Legendre transforms:} Each thermodynamic potential has a differential that pairs conjugate variables. The exactness condition for each gives the corresponding Maxwell relation.
  \item \textbf{From the cyclic relation:} The triple product rule $(\partial x/\partial y)_{z}(\partial y/\partial z)_{x}(\partial z/\partial x)_{y} = -1$ combined with Maxwell relations can derive additional identities.
\end{enumerate}
\section*{Verifications}

Maxwell relations are exact consequences of the laws of thermodynamics and are verified indirectly through their predictions. For example, for an ideal gas, the prediction that $(\partial S/\partial V)_{T} = nR/V$ is confirmed by calorimetric measurements of entropy. The adiabatic relation $\gamma = C_{P}/C_{V}$ and its derived expressions for sound speed are experimentally consistent with Maxwell relation predictions.
\section*{Applications}

\begin{itemize}
  \item \textbf{Equation of state:} For an ideal gas ($PV = nRT$), Maxwell relation~\eqref{eq:maxwell3} gives $(\partial S/\partial V)_{T} = nR/V$, which upon integration gives the Sackur--Tetrode entropy.
  \item \textbf{Specific heats:} The relation between $C_{P}$ and $C_{V}$ involves a Maxwell relation to evaluate the remaining derivative.
  \item \textbf{Thermal expansion:} The coefficient of thermal expansion $\alpha = (1/V)(\partial V/\partial T)_{P}$ enters Maxwell relation~\eqref{eq:maxwell4}.
  \item \textbf{Clausius--Clapeyron equation:} Derivation of the slope of phase boundaries in $P$--$T$ diagrams uses a Maxwell relation.
\end{itemize}
\section*{What Richard Feynman Would Have Asked}

\subsection*{Plain-English Translation}
\textit{``What is the Maxwell relation really saying, if I describe it without any calculus?''}

It says that nature is consistent: if you push on a system in two different ways---say, squeeze it and heat it---the way temperature responds to squeezing is locked in by the way pressure responds to heating. You cannot change one without the other changing in a predictable, complementary way. The four Maxwell relations are like four different ``consistency checks'' for the bookkeeping of energy in thermodynamics.

\subsection*{Localized Mechanics}
\textit{``At the level of the mathematics, why do these relations hold, and what is the essential structure they encode?''}

The key mathematical fact is Schwarz's theorem: if a function $F(x,y)$ has continuous second partial derivatives, then $\partial^{2}F/\partial x\,\partial y = \partial^{2}F/\partial y\,\partial x$. In thermodynamics, the differential $dU = T\,dS - P\,dV$ has $T$ and $-P$ as partial derivatives of $U$ with respect to $S$ and $V$, respectively. Schwarz's theorem then forces $\partial T/\partial V = -\partial P/\partial S$ (with the appropriate variables held constant). This is not a physical law but a mathematical constraint: the state function $U$ is a smooth function, and smooth functions have order-independent mixed derivatives. The physics enters through the identification of the conjugate pairs $(T, S)$ and $(P, V)$, but the relation itself is pure calculus.

\subsection*{Testing Extreme Limits}
\textit{``Do the Maxwell relations still hold in extreme conditions---near absolute zero, near a critical point, or in extreme pressures?''}

Near absolute zero, the third law of thermodynamics ($S \to 0$ as $T \to 0$) ensures that the potentials remain smooth, so Maxwell relations hold---though the derivatives may approach singular values (e.g., $C_{P}/T \to$ finite as $T \to 0$ for metals). Near a critical point, the response functions ($C_{V}$, $\kappa_{T}$) diverge, and the Maxwell relations predict specific relationships between these divergences---this is the basis of critical exponent relations in scaling theory. In extreme pressures (e.g., neutron star interiors), the relations still hold if the matter is in local thermodynamic equilibrium, even if the equation of state is highly nonlinear. The Maxwell relations are remarkably robust.

\subsection*{Dimensionality and Geometry}
\textit{``Is there a geometric picture of the Maxwell relations, like a surface or a manifold?''}

Yes. The thermodynamic potentials define surfaces in the space of state variables. For example, $U(S,V)$ is a surface in three-dimensional $(U, S, V)$ space. The partial derivatives $\partial U/\partial S = T$ and $\partial U/\partial V = -P$ are the slopes of this surface. The Maxwell relation $\partial T/\partial V = -\partial P/\partial S$ is simply the statement that the surface is smooth (no ``kinks'' except at phase transitions). The Legendre transforms that connect $U$, $H$, $F$, and $G$ are like changing coordinates on this surface, projecting it onto different planes. The four Maxwell relations are the integrability conditions for the surface in each of these coordinate systems.

\subsection*{Cross-Disciplinary Connection}
\textit{``Where does this structure---exact differentials giving constraint equations---appear outside of thermodynamics?''}

In differential geometry, the condition $d^{2} = 0$ (the exterior derivative squares to zero) underlies the Poincaré lemma and de Rham cohomology; Maxwell's equations themselves are an expression of $dF = 0$ and $d{\star}F = J$ in the language of differential forms. In fluid mechanics, the Crocco theorem relates vorticity to entropy gradients through a constraint that is structurally identical to a Maxwell relation. In economics, the comparative-statics identities (Shephard's lemma, Roy's identity) are derived from the exactness of utility and expenditure functions---a direct analogue. The deep lesson is that whenever a system is described by a smooth potential with conjugate variable pairs, constraint equations like the Maxwell relations must hold.

% ============================================================
% CHAPTER 41 — Maxwell–Boltzmann Distribution
% ============================================================
\section*{FAQ}

\begin{enumerate}
\item \textbf{Are the Maxwell relations specific to ideal gases?}
No. They hold for any system described by well-defined thermodynamic potentials. Their application to an ideal gas is simply the most common textbook example.

\item \textbf{Why are there exactly four?}
There are four because there are four standard thermodynamic potentials ($U$, $H$, $F$, $G$), each with a natural pair of independent variables, and each exact differential gives one relation.

\item \textbf{What is the physical meaning of $(\partial T/\partial V)_{S} = -(\partial P/\partial S)_{V}$?}
It says that if you expand a gas adiabatically (no heat exchange), the temperature drop rate is determined by the same physics as how the pressure responds to entropy changes at constant volume. Both sides measure the same thermodynamic response from different perspectives.
\end{enumerate}
\section*{Important Bibliography}

\begin{enumerate}
\item Callen, H.B., \textit{Thermodynamics and an Introduction to Thermostatistics}, 2nd ed. (Wiley, 1985), Chapter 5.
\item Pathria, R.K. and Beale, P.D., \textit{Statistical Mechanics}, 4th ed. (Academic Press, 2021), Chapter 3.
\item Schroeder, D.V., \textit{An Introduction to Thermal Physics} (Addison-Wesley, 2000), Chapter 5.
\item Reif, F., \textit{Fundamentals of Statistical and Thermal Physics} (McGraw-Hill, 1965), Chapter 3.
\end{enumerate}

